\documentclass{article}

\usepackage{lipsum}
\usepackage[margin=4cm, includefoot]{geometry}
\usepackage{graphicx}
\usepackage{float}
\usepackage{hyperref}

% Header and footer
\usepackage{fancyhdr}
\pagestyle{fancy}

\rhead{}
\lhead{}
\fancyfoot{}
\fancyfoot[R]{\thepage}
\renewcommand{\headrulewidth}{0pt}
\renewcommand{\footrulewidth}{0pt}
%

\begin{document}

\begin{titlepage}
	\begin{center}
		\line(1,0){400}\\
		[6mm]
		\huge{\bfseries Architectural and Functional Requirements}\\
		\line(1,0){400}\\
			[5mm]
			\includegraphics[width=150px]{images/AWorldOfPlants.png}
			\\
		[5mm]
		\large\textbf{Project:}\\A World of Things\\
		[3mm]
		\large\textbf{Client:}\\Julian Hambleton-Jones\\
		[3mm]
		\large \textbf{Team:}\\Funge\\
		\line(1,0){400}\\
		[5mm]
		\large \textbf{Team Members:}\\
		[3mm]
		\large 14214742 - Matthew Botha\\
		\large 14446619 - Gian Paolo Buffo\\
		\large 14027021 - Matthias Harvey\\
        \large 14035538 - Dillon Heins\\[3mm]
	\end{center}
\end{titlepage}

\cleardoublepage
\thispagestyle{empty}
\tableofcontents
\cleardoublepage
\setcounter{page}{1}

\section{Introduction}
	The Internet of Things (IoT) is a development of the Internet which involves the networking of every day physical devices allowing them to send and receive data. These devices are embedded with electronics, sensors, software as well as some form of Internet or network connectivity.


	IoT is a relatively new development and has a large possibility of becoming a ubiquitous technology as well as allowing us to view the world from a different perspective. The potential it contains for innovation is endless.


	We as a group were given the opportunity to use the Amazon Web Services (AWS) IoT platform to create an IoT project of our own desires. This document details our 'A World of Things' project.

\section{Vision}
	The aim of this project is to build an innovative Internet of Things solution through the use of the Amazon Web Services Internet of Things hardware platform as well as the Amazon Web Services Cloud.
	
	No problem was formally defined and hence it was up to us to determine what solution we would be creating. We opted to create a solution which focuses on the education of students with regard to agriculture and plant sciences. We plan to motivate students to become interested in agriculture by creating an "Internet of Plants".
	
	This involves the networking and monitoring of living plants in order to analyse aspects of their environment - such as water intake, lighting, humidity, moisture, and temperature. It also involves the controlling of the environment of the plants through the adjustment of the amount of water they receive, the wavelengths of their lights and air flow through the use of fans.
	
	Users should be able to view data gathered about their plants through the use of graphs so that they can make informed decisions about what adjustments to the environment they should make.
	
	The main purpose of the platform is to encourage younger generations to become interested in agriculture and, in doing so, help stimulate South Africa's agricultural industry. We plan to implement gamification on our platform as a way to encourage users to participate.
	
	Given enough time, AI learning could be employed to optimise the conditions under which plants grow. By combining automation and AI learning, we could create an interesting challenge for the users: Grow a plant better than our control plant grown with the help of AI.

\section{Background}
	This project is part of a Software Engineering module (COS 301) for the University of Pretoria. Our client is Julian Hambleton-Jones (AWS) and our project was to create anything we wanted using the AWS IoT platform. What we decided upon was the creation of a network of living plants which are to be analysed and looked after through the use of IoT devices and our software system.
	
\subsection{Future Business/Research Opportunities}
	
\subsection{The Client's Problem}

\section{Important Terminology}
\begin{itemize}
	\item \textbf{A World of Plants} Name of the entire system.
	\item \textbf{AWS} Amazon Web Services. A secure cloud services platform offering compute power, database storage, content delivery and other functionality.
	\item \textbf{Plant Box} The enclosure which contains the plants, with all the monitoring and automation equipment.
	\item \textbf{Achievements} Gamified awards given to users for completing certain tasks.
\end{itemize}

\section{Architecture Requirements}

\subsection{Access Channel Requirements}

\subsubsection{Human Access Channels}
Users will be able to register and log into a Web interface. From the Web interface, they can associate their own Plant Boxes, monitor conditions such as temperature and humidity, control the Plant Box conditions and view their Achievements. The Web interface should be accessible from any modern browser, such as Google Chrome, Mozilla Firefox, Microsoft Edge, Apple Safari and latest Internet Explorer. The Web interface will be designed under strict UX, UI and World Wide Web Consortium (W3C) standards. This includes a responsive design that will work on any screen size. Plant Boxes do not need to be connected to the Internet at all times - the Web interface will store the latest settings and sync them with the Plant Boxes when a connection can be established.     

\subsubsection{System Access Channels}
AWS API Gateway will be used expose the system's API and map it to the serverless backend provided by AWS Lambda. MQTT messages will be used to communicate between the system backend and the individual Plant Boxes. MQTT is an ISO standard (ISO/IEC PRF 20922) publish-subscribe-based "lightweight" messaging protocol for use on top of the TCP/IP protocol. It is often used to facilitate IoT communications.

\subsection{Quality Requirements}

\subsubsection{Performance}

\subsubsection{Usability}

\subsubsection{Reliability}

\subsubsection{Scalability}

\subsubsection{Flexibility}

\subsubsection{Security}

\subsubsection{Auditability/Monitorability}

\subsubsection{Integrability}

\subsection{Integration Requirements}

\subsection{Architecture Constraints}

\section{Architecture Design}

\subsection{Infrastructure}

\subsection{Services}

\subsubsection{TODO add services here}

\subsection{Tactics}

\subsubsection{Flexibility}

\subsubsection{Reliability}

\subsubsection{Security}

\subsubsection{Auditability}

\section{Database and Persistence}

\section{Process Specification}

\subsection{Server}

\subsection{Frontend}

\section{Functional Requirements}

\subsection{Use Case Prioritisation}

\subsection{Use Case/Service Contracts}

\subsubsection{Add use cases/service contracts}

\section{Technologies}

\subsection{Todo add technologies like framework, web container etc}

\section{Initial Design}

\section{Open Issues}

\end{document}
